\subsection{黄金半径链、Riccati 固定点与端点缺陷证书的统一极限结构}\label{subsec:conclusion-golden-radius-riccati-logistic-endpoint-interface}
沿用推论 \ref{cor:pom-Lk-area-law-qform}、推论 \ref{cor:pom-Lk-boundary-fixed-point-riccati}、推论 \ref{cor:pom-Lk-surface-free-energy}、命题 \ref{prop:pom-Lk-boundary-riccati-recursion} 与定理 \ref{thm:xi-rh-iff-uniform-comoving-defect} 的记号。
本小节将黄金尺度、边界 Riccati 递推与 Jensen 缺陷接口压缩为一条统一的可审计定理链。

\begin{theorem}[黄金半径链的 Fibonacci 中和极限]\label{thm:conclusion-golden-radius-fibonacci-neutralization-limit}
\leanverified{paper\_conclusion\_golden\_radius\_fibonacci\_neutralization\_limit}
设
\[
\varphi:=\frac{1+\sqrt5}{2},
\qquad
F_0=0,\ F_1=1,\ F_{n+2}=F_{n+1}+F_n,
\]
并定义黄金半径链
\[
\rho_m:=\tanh\!\Bigl(\frac{m\log\varphi}{2}\Bigr),\qquad m\ge 1.
\]
对任意固定整数 \(k\)，有极限
\[
\lim_{m\to\infty}\rho_m^{\,F_{m+k}}
=
\exp\!\Bigl(-\frac{2}{\sqrt5}\,\varphi^k\Bigr).
\]
\end{theorem}
\begin{proof}
由
\[
\rho_m=\frac{\varphi^m-1}{\varphi^m+1}
\]
得
\[
1-\rho_m=\frac{2}{\varphi^m+1}=2\varphi^{-m}+O(\varphi^{-2m}),
\]
故
\[
\log \rho_m
=\log\!\bigl(1-(1-\rho_m)\bigr)
=-2\varphi^{-m}+O(\varphi^{-2m}).
\]
另一方面，Binet 公式给出
\[
F_{m+k}=\frac{\varphi^{m+k}}{\sqrt5}+O(\varphi^{-(m+k)}).
\]
两式相乘得
\[
F_{m+k}\log\rho_m
=-\frac{2}{\sqrt5}\,\varphi^k+o(1),
\]
指数化即得结论。证毕。
\end{proof}

\begin{corollary}[中和常数的乘法 Fibonacci 递推]\label{cor:conclusion-golden-radius-neutralization-multiplicative-recursion}
\leanverified{paper\_conclusion\_golden\_radius\_neutralization\_multiplicative\_recursion}
定义
\[
r_k:=\exp\!\Bigl(-\frac{2}{\sqrt5}\,\varphi^k\Bigr),\qquad k\in\ZZ.
\]
则
\[
r_{k+2}=r_{k+1}\,r_k,
\qquad
r_{k+1}=r_k^{\,\varphi}.
\]
\end{corollary}
\begin{proof}
由 \(\varphi^2=\varphi+1\)，有
\[
r_{k+1}r_k
=\exp\!\Bigl(-\frac{2}{\sqrt5}(\varphi^{k+1}+\varphi^k)\Bigr)
=\exp\!\Bigl(-\frac{2}{\sqrt5}\varphi^{k+2}\Bigr)
=r_{k+2}.
\]
同理
\[
r_k^{\,\varphi}
=\exp\!\Bigl(-\frac{2}{\sqrt5}\varphi^{k+1}\Bigr)
=r_{k+1}.
\]
证毕。
\end{proof}

\begin{theorem}[Riccati 固定点与双 Cayley 坐标的精确同构]\label{thm:conclusion-riccati-double-cayley-exact-isomorphism}
\leanverified{paper\_conclusion\_riccati\_double\_cayley\_exact\_isomorphism}
设 \(t>0\)，并记
\[
\eta(t):=\operatorname{arsinh}\!\Bigl(\frac{\sqrt t}{2}\Bigr),
\qquad
q(t):=e^{-2\eta(t)}.
\]
由推论 \ref{cor:pom-Lk-boundary-fixed-point-riccati}，\(q(t)\) 为
\[
q^2-(t+2)q+1=0
\]
的小根。定义
\[
\rho(t):=\tanh\!\Bigl(\frac{\eta(t)}{2}\Bigr).
\]
则
\[
q(t)=\Bigl(\frac{1-\rho(t)}{1+\rho(t)}\Bigr)^2,
\qquad
t=\frac{(1-q(t))^2}{q(t)}.
\]
\end{theorem}
\begin{proof}
由恒等式
\[
\tanh\!\Bigl(\frac{x}{2}\Bigr)=\frac{1-e^{-x}}{1+e^{-x}}
\]
可得
\[
e^{-x}=\frac{1-\tanh(x/2)}{1+\tanh(x/2)}.
\]
取 \(x=\eta(t)\) 即
\[
e^{-\eta(t)}=\frac{1-\rho(t)}{1+\rho(t)},
\]
故
\[
q(t)=e^{-2\eta(t)}
=\Bigl(\frac{1-\rho(t)}{1+\rho(t)}\Bigr)^2.
\]
由 \(q^2-(t+2)q+1=0\) 直接整理得
\[
t=\frac{q^2-2q+1}{q}=\frac{(1-q)^2}{q}.
\]
证毕。
\end{proof}

\begin{theorem}[二次数域的分歧素数与最小素数账本]\label{thm:conclusion-quadratic-field-min-prime-ledger-discriminant}
设 \(K/\QQ\) 为二次数域，\(\mathcal O_K\) 为其整数环，并记其基本判别式为 \(\Disc(K)\)。
则二次数域的分歧素数集合满足
\[
\mathrm{Ram}(K)=\{\,p\in\mathbb P:\ p\mid \Disc(K)\,\}.
\]
在“局部提升一致 + 实现可逆”之记录轴口径下，任何试图在所有良素数方向上统一实现 \(\mathcal O_K\) 的同余塔一致读出/提升的方案，其素数账本支撑必包含 \(\mathrm{Ram}(K)\)，并且在二次数域情形下该必要条件亦为最小条件。
\end{theorem}
\begin{proof}
等式 \(\mathrm{Ram}(K)=\{p:\ p\mid \Disc(K)\}\) 为代数数论的标准判别（例如由局部化 \(\mathcal O_K\otimes\ZZ_p\) 的 \'etale/非 \'etale 性刻画）。
\par
当 \(p\nmid \Disc(K)\) 时，\(\mathcal O_K\otimes\ZZ_p\) 为秩 \(2\) 的有限 \'etale \(\ZZ_p\)-代数；因此模 \(p\) 的分解型在 Hensel 引理下具有唯一提升，从而局部同余塔可由有限层数据一致控制。
当 \(p\mid \Disc(K)\) 时，\(\mathcal O_K\otimes\ZZ_p\) 非 \'etale，模 \(p\) 约化出现重根/非约化结构，导致同余塔上出现无法由有限层统一压缩的提升分叉；因此任何在该方向上要求可逆且提升一致的实现都必须显式保留该素数方向。
二次数域中分歧仅由上述两类局部情形决定，因而最小素数账本支撑恰为 \(\mathrm{Ram}(K)\)。证毕。
\end{proof}

\begin{corollary}[Riccati 二次域参数的分歧素数支撑]\label{cor:conclusion-riccati-quadratic-field-ramification-ledger}
在定理 \ref{thm:conclusion-riccati-double-cayley-exact-isomorphism} 的记号下，取 \(t\in\QQ_{>0}\) 并令
\[
K_t:=\QQ\!\left(\sqrt{t(4+t)}\right).
\]
则 \(K_t/\QQ\) 的最小素数账本支撑恰为其分歧素数集合，亦即为 \(\Disc(K_t)\) 的素支撑集合。
\end{corollary}
\begin{proof}
由 \(K_t=\QQ(\sqrt{t(4+t)})\) 可知其为二次数域，直接应用定理 \ref{thm:conclusion-quadratic-field-min-prime-ledger-discriminant} 即得。证毕。
\end{proof}

\endinput



\begin{theorem}[黄金耦合的有限尺寸刚性与参数反演界]\label{thm:conclusion-golden-coupling-finite-k-rigidity}
定义
\[
\Lambda_k(t):=\frac1k\log\det(L_k+tI),\qquad t>0,\ k\ge 1.
\]
由推论 \ref{cor:pom-Lk-area-law-qform}，
\[
\log\det(L_k+tI)
=2k\,\eta(t)-\log(1+q(t))+\log\!\bigl(1+q(t)^{2k+1}\bigr).
\]
因此对所有 \(t>0\) 与 \(k\ge 1\) 有一致界
\[
\bigl|\Lambda_k(t)-2\eta(t)\bigr|\le \frac{\log 2}{k}.
\]
进一步，若
\[
\bigl|\Lambda_k(t)-2\log\varphi\bigr|\le \delta,
\]
则
\[
|\eta(t)-\log\varphi|
\le
\frac{\delta}{2}+\frac{\log 2}{2k},
\qquad
t=4\sinh^2(\eta(t)).
\]
并且
\[
\frac{\dd t}{\dd\eta}\Big|_{\eta=\log\varphi}=2\sqrt5.
\]
\end{theorem}
\begin{proof}
由 \(0<q(t)<1\) 知
\[
0<\log(1+q(t)^{2k+1})<\log 2,\qquad
0<\log(1+q(t))<\log 2.
\]
于是
\[
\bigl|-\log(1+q(t))+\log(1+q(t)^{2k+1})\bigr|\le \log 2,
\]
除以 \(k\) 得第一不等式。

再由三角不等式，
\[
|\eta(t)-\log\varphi|
\le
\frac12|\Lambda_k(t)-2\eta(t)|+\frac12|\Lambda_k(t)-2\log\varphi|
\le
\frac{\log 2}{2k}+\frac{\delta}{2}.
\]
由 \(\eta(t)=\operatorname{arsinh}(\sqrt t/2)\) 得 \(t=4\sinh^2\eta(t)\)。
最后
\[
\frac{\dd t}{\dd\eta}=8\sinh\eta\cosh\eta,
\]
在 \(\eta=\log\varphi\) 处，\(\sinh(\log\varphi)=\frac12\)、\(\cosh(\log\varphi)=\frac{\sqrt5}{2}\)，故导数为 \(2\sqrt5\)。证毕。
\end{proof}

\begin{theorem}[表面项的 Logistic 输运结构]\label{thm:conclusion-surface-term-logistic-transport}
\leanverified{paper\_conclusion\_surface\_term\_logistic\_transport}
设
\[
f(t):=2\eta(t),\qquad
g(t):=-\log(1+q(t)),
\qquad
\sigma(t):=\frac{q(t)}{1+q(t)}\in\Bigl(0,\frac12\Bigr).
\]
则成立：
\begin{enumerate}
  \item 以 \(f\) 为演化参数的 Logistic 动力学
  \[
  \frac{\dd \sigma}{\dd f}=-\sigma(1-\sigma),
  \qquad
  \sigma(0)=\frac12;
  \]
  \item 显式解
  \[
  \sigma(t)=\frac{1}{1+e^{f(t)}},
  \qquad
  g(t)=-\log\!\bigl(1+e^{-f(t)}\bigr).
  \]
\end{enumerate}
\end{theorem}
\begin{proof}
由 \(q(t)=e^{-f(t)}\) 得
\[
q'(t)=-q(t)f'(t).
\]
于是
\[
\frac{\dd \sigma}{\dd t}
=\frac{q'(t)}{(1+q(t))^2}
=-\frac{q(t)}{(1+q(t))^2}f'(t)
=-\sigma(t)(1-\sigma(t))f'(t),
\]
故
\[
\frac{\dd\sigma}{\dd f}=-\sigma(1-\sigma).
\]
当 \(t\downarrow 0\) 时，\(\eta(t)\downarrow0\)、\(q(t)\uparrow1\)，从而 \(\sigma(0)=1/2\)。

解 Logistic 方程：
\[
\log\frac{\sigma}{1-\sigma}=-f+C,
\]
由初值 \(C=0\)，得
\[
\sigma=\frac{1}{1+e^f}.
\]
再由 \(g=-\log(1+q)\) 与 \(q=e^{-f}\) 即得
\[
g=-\log(1+e^{-f}).
\]
证毕。
\end{proof}

\begin{theorem}[边界反射误差与有限尺寸校正的导数同一性]\label{thm:conclusion-boundary-reflection-derivative-identity}
\leanverified{paper\_conclusion\_boundary\_reflection\_derivative\_identity}
定义
\[
s_k(t):=\log\!\bigl(1+q(t)^{2k+1}\bigr),
\qquad
\Delta_k(t):=q_k(t)-q(t),
\]
其中 \(q_k(t)\) 与 \(q(t)\) 如命题 \ref{prop:pom-Lk-boundary-riccati-recursion}，并满足
\[
\Delta_k(t)=\frac{(1-q(t)^2)\,q(t)^{2k}}{1+q(t)^{2k+1}}>0.
\]
则对任意 \(t>0\)、\(k\ge 0\) 有精确恒等式
\[
\Delta_k(t)
=
-\,\frac{1-q(t)^2}{(2k+1)\,q(t)\,f'(t)}
\cdot\frac{\dd}{\dd t}s_k(t),
\qquad
f'(t)=\frac{1}{\sqrt{t(4+t)}}.
\]
\end{theorem}
\begin{proof}
由链式法则，
\[
\frac{\dd}{\dd t}s_k(t)
=\frac{(2k+1)q(t)^{2k}q'(t)}{1+q(t)^{2k+1}}.
\]
又 \(q'(t)=-q(t)f'(t)\)，故
\[
\frac{\dd}{\dd t}s_k(t)
=-(2k+1)\frac{q(t)^{2k+1}f'(t)}{1+q(t)^{2k+1}}.
\]
与
\[
\Delta_k(t)=\frac{(1-q(t)^2)\,q(t)^{2k}}{1+q(t)^{2k+1}}
\]
比较即得结论。证毕。
\end{proof}

\begin{corollary}[黄金点的最小窗长取整律]\label{cor:conclusion-boundary-reflection-min-window-length-integer-lock}
\leanverified{paper\_conclusion\_boundary\_reflection\_min\_window\_length\_integer\_lock}
沿用定理 \ref{thm:conclusion-boundary-reflection-derivative-identity} 的记号。
对任意 \(t>0\) 与容差 \(\varepsilon\in(0,1-q(t)^2)\)，定义最小窗长
\[
k_{\min}(t,\varepsilon):=\min\{k\in\ZZ_{\ge 0}:\ \Delta_k(t)\le \varepsilon\}.
\]
则 \(k_{\min}(t,\varepsilon)\) 满足双边估计
\[
\left\lceil\frac{\log\!\bigl((1-q(t)^2)/\varepsilon\bigr)}{2\log(1/q(t))}\right\rceil-1
\le
k_{\min}(t,\varepsilon)
\le
\left\lceil\frac{\log\!\bigl((1-q(t)^2)/\varepsilon\bigr)}{2\log(1/q(t))}\right\rceil.
\]
并且在 \(t=1\)（从而 \(q(1)=\varphi^{-2}\)）时，对任意整数 \(m\ge 1\) 有严格恒等式
\[
k_{\min}\bigl(1,\varphi^{-4m}\bigr)=m.
\]
\end{corollary}
\begin{proof}
由
\(
0<\Delta_k(t)=\frac{(1-q^2)q^{2k}}{1+q^{2k+1}}
\)
可知 \(\Delta_k(t)\) 关于 \(k\) 单调递减。
并且对所有 \(k\ge 0\) 有
\[
\frac{1-q(t)^2}{1+q(t)}\,q(t)^{2k}\le \Delta_k(t)\le (1-q(t)^2)\,q(t)^{2k}.
\]
由此可将条件 \(\Delta_k(t)\le \varepsilon\) 夹逼为关于 \(q(t)^{2k}\) 的对数不等式，从而得到所示双边估计。

在 \(t=1\) 时，\(q(1)\) 为方程 \(q^2-3q+1=0\) 的小根，即 \(q(1)=\varphi^{-2}\)。
记 \(q:=\varphi^{-2}\)。
由显式尾差式
\(
\Delta_k(1)=\frac{(1-q^2)q^{2k}}{1+q^{2k+1}}
\)，
对 \(k=m\) 有 \(\Delta_m(1)<q^{2m}=\varphi^{-4m}\)，故 \(k_{\min}\le m\)。
另一方面，对 \(k=m-1\) 有
\[
\Delta_{m-1}(1)
=\frac{(1-q^2)q^{2m-2}}{1+q^{2m-1}}
>
\frac{(1-q^2)q^{2m-2}}{1+q}
=\varphi^{-4m}\cdot\frac{\varphi^{4}-1}{1+\varphi^{-2}}
>\varphi^{-4m},
\]
从而 \(k_{\min}\ge m\)。
两式合并即得 \(k_{\min}(1,\varphi^{-4m})=m\)。证毕。
\end{proof}

\begin{theorem}[黄金--Fibonacci 混合幂的超指数衰减]\label{thm:conclusion-golden-fibonacci-mixed-power-superexponential}
沿用定理 \ref{thm:conclusion-golden-radius-fibonacci-neutralization-limit} 的 \(\rho_m\) 与 \(F_m\)，并定义
\[
R_m:=\rho_m^{\,F_m}.
\]
则
\[
\lim_{m\to\infty}R_m=\exp\!\Bigl(-\frac{2}{\sqrt5}\Bigr),
\]
且
\[
R_m^{\,F_{m+2}}
=
\exp\!\Bigl(-\frac{2}{5}\,\varphi^{m+2}+o(\varphi^m)\Bigr),
\qquad m\to\infty.
\]
因此 \(R_m^{F_{m+2}}\) 以 \(\exp(-c\varphi^m)\) 级别趋于 \(0\)，严格快于任意 \(e^{-cm}\) 型指数衰减。
\end{theorem}
\begin{proof}
第一式是定理 \ref{thm:conclusion-golden-radius-fibonacci-neutralization-limit} 在 \(k=0\) 的特例。

第二式由
\[
\log\!\bigl(R_m^{F_{m+2}}\bigr)=F_{m+2}F_m\log\rho_m
\]
与渐近
\[
F_n=\frac{\varphi^n}{\sqrt5}+O(\varphi^{-n}),
\qquad
\log\rho_m=-2\varphi^{-m}+O(\varphi^{-2m})
\]
相乘得到
\[
F_{m+2}F_m\log\rho_m
=-\frac{2}{5}\varphi^{m+2}+o(\varphi^m).
\]
指数化即得结论。证毕。
\end{proof}

\begin{conjecture}[端点窗口的 Logistic 普适律与 RH 审计接口]\label{conj:conclusion-endpoint-logistic-universality-rh-interface}
记共动圆盘化对象与 Jensen 缺陷
\[
F_{\zeta,x_0}(w):=\Xi_\zeta\!\bigl(s_{x_0}(w)\bigr),\qquad
D_{\zeta,x_0}(\varrho)
\]
如 \S\ref{subsubsec:xi-comoving-defect-delta-bound} 所定义。
存在一条与 \(x_0\) 一致的端点窗口结构，使下述三层机制同时成立：
\begin{enumerate}
  \item 存在尺度参数 \(t=t(\varrho)\to\infty\)（\(\varrho\uparrow1\)）使主导输运仅经 \(f(t)\) 与 \(q(t)=e^{-f(t)}\) 进入，且满足与定理 \ref{thm:conclusion-surface-term-logistic-transport} 同型的 Logistic 演化；
  \item 存在与命题 \ref{prop:pom-Lk-boundary-riccati-recursion} 同型的有限窗反射对象 \(q_k(t)\)，使
  \[
  \sup_{x_0\in\RR}D_{\zeta,x_0}(\varrho)
  \le
  \mathcal E_{\mathrm{bulk}}(t)+C\bigl(q_k(t)-q(t)\bigr)+o(1),
  \]
  且 \(\mathcal E_{\mathrm{bulk}}(t)\) 具有显式高阶衰减；
  \item 在黄金分辨率链
  \[
  \rho_m=\tanh\!\Bigl(\frac{m\log\varphi}{2}\Bigr)
  \]
  下，存在线性耦合 \(k\sim m\) 使右端收敛到可触发定理 \ref{thm:xi-rh-iff-uniform-comoving-defect} 与定理 \ref{thm:xi-jensen-single-radius-band-exclusion} 的统一缺陷门槛。
\end{enumerate}
若该猜想成立，则 RH 判据可由“Logistic--Riccati 有限维证书链 + Jensen 缺陷上界”实现端点化审计闭合。
\end{conjecture}

\begin{remark}[统一读出链]\label{rem:conclusion-golden-radius-riccati-logistic-chain}
上述结构将
\[
\text{黄金半径端点几何}
\Longrightarrow
\text{Riccati 固定点与有限窗误差}
\Longrightarrow
\text{Logistic 输运与导数同一性}
\Longrightarrow
\text{Jensen 缺陷的 RH 审计接口}
\]
压缩为同一条可验证的参数化链路。
\end{remark}
